
En las últimas actualizaciones de \textit{RTXI}, se incluye una nueva funcionalidad, que se realiza de forma inevitable cada vez que se hace una conexión, la función `find\_cycle`. Esta función busca si hay un ciclo en las conexiones que se realizan en \textit{RTXI} entre plugins, y en caso de que lo haya, se evita que se llegue a crear (Código \ref{COD:CHECKCYCLE}). 

Esto supone un problema para usar las herramientas empleadas hasta ahora, y para crear cualquier circuito de ciclo cerrado con esta herramienta. Por ello se han aportado dos soluciones, una es una versión con dicha directiva eliminada, y la otra es la modificación de la estructura de las conexiones de los "bloques" (así es como se refiere a los registros I/O \textit{RTXI}), creando una variable booleana llamada `check\_cycle` que indica que se debe revisar si la conexión ocasiona un ciclo. Por otro lado, se ha añadido un botón al componente `Connector`, llamado `Check Cycle`, este representa el valor para comprobar la conexión, de manera que si este está desactivado, no se buscarán ciclos (Repositorio \ref{SEC:RTXIPREEMPTREPO}).

No obstante, la modificación con el botón no está recomendada en la versión actual, dado que la función que se usa para buscar un ciclo es recursiva, por lo que si se desactiva la variable, se crea un ciclo y luego se intenta crear otro con dicha variable activa, se entra en una recursión infinita y el programa termina con `Segmentation fault`. Esto no pasa en la versión en la que la comprobación ha sido desactivada, dado que no se busca en ninguna conexión si se ha creado un ciclo. Sin embargo, hay otro problema mayor ligado a los ciclos cerrados en \textit{RTXI}, y es que es importante realizar las conexiones de cualquier otro componente base de la aplicación antes de realizar la conexión en ciclo entre los plugins, es decir, si se desea visualizar una frecuencia con el \textit{Oscilloscope}, recoger datos con el \textit{Data Recorder} o mostrar datos con el \textit{RT Benchmarks}, se deben de instanciar estos primero, pues si durante el ciclo cerrado se interactúa con dichos componentes, los plugins conectados al ciclo quedarán inutilizables, obligando al usuario a o bien eliminarlos y cargarlos de nuevo, o bien a liberar las conexiones relacionadas en el ciclo cerrado y las realizadas en el \textit{Oscilloscope} y el \textit{Data Recorder}, en el caso de \textit{RT Benchmarks} basta con cerrarlo.

Estos problemas nacen de la manera en la que \textit{RTXI} concibe las conexiones con los plugins, pues durante su ejecución, realiza una búsqueda topológica basándose en las conexiones, y esto se hace usando  \textit{Directed Acyclic Graphs} (\textit{DAGs}), especificado por el propio mantenedor principal del proyecto en Github \cite{IvanFValerio2024}. Como solución a esto, se puede, o bien usar otro ordenador para instanciar ahí la neurona modelo, o bien usar los propios canales de la DAQ para transmitir la información entre los dos plugins cuya conexión generaría el ciclo, dado que este ciclo cerrado solo se genera si no hay un dispositivo entre ambos plugins (Sección \ref{SS:HRHRELECTRO}).

\CPPCode[COD:CHECKCYCLE]{Comprueba si hay un ciclo.}{En esta figura se muestra parte del código de la función `connect`, de \textit{RTXI}.}{rtxi/src/rt.cpp}{65}{68}{1}

